../cictikz/src/cictikz/data/tex/cictikz_preamble.tex

% Currents always run in loops. Left: the loop you drew - a source, a
% wire, a load, and the return. Right: the loop you get when you forget
% the return - the current still closes, through whatever shared ground
% is available, and everything on that path sees it.

\begin{circuitikz}[american, thick, transform shape, circuit ee IEC]

  ../cictikz/src/cictikz/data/tex/cictikz_lib.tex

  % ---------------- (a) the loop you drew ----------------
  \draw (0,0) to [V] (0,2.4);
  \draw (0,2.4) -- (3.6,2.4);
  \draw (3.6,0.6) \vresistor{$R$};
  \draw (3.6,2.4) -- (3.6,2.2);
  \draw (3.6,0.6) -- (3.6,0);
  \draw (3.6,0) -- (0,0);
  % the current, one loop
  \draw[blue, very thick, ->] (0.7,2.05) -- (2.9,2.05);
  \node[blue, anchor=south] at (1.8,2.1) {$I$};
  \draw[blue, very thick, ->] (2.9,0.35) -- (0.7,0.35);
  \node[anchor=north, align=center] at (1.8,-0.7) {(a) the loop you drew};

  % ---------------- (b) the loop you got ----------------
  \begin{scope}[shift={(7.2,0)}]
    \draw (0,0) to [V] (0,2.4);
    \draw (0,2.4) -- (3.6,2.4);
    \draw (3.6,0.6) \vresistor{$R$};
    \draw (3.6,2.4) -- (3.6,2.2);
    \draw (3.6,0.6) -- (3.6,0);
    % no tidy return: the current finds its own way home
    \draw (3.6,0) -- (4.3,0) -- (4.3,-1.2) -- (-0.7,-1.2) -- (-0.7,0) -- (0,0);
    \node[anchor=west, scale=0.8] at (4.4,-0.6) {shared ground};
    \draw[blue, very thick, ->] (0.7,2.05) -- (2.9,2.05);
    \node[blue, anchor=south] at (1.8,2.1) {$I$};
    \draw[red, very thick, ->] (3.95,-0.85) -- (-0.35,-0.85);
    \node[red, anchor=north] at (1.8,-1.35) {the return you forgot still exists};
    \node[anchor=north, align=center] at (1.8,-2.15) {(b) the loop you got};
  \end{scope}

\end{circuitikz}
\end{document}
